\section{Results Analysis}


This section describes the planned analysis methods for evaluating the impact of
different prompt engineering strategies.


The analysis focuses on comparing whether structured prompting improves
algorithmic problem-solving performance.



\subsection{Overall Performance Comparison}


The first analysis compares the overall performance of the three experimental
groups:


\begin{itemize}

\item Baseline Prompt.

\item Optimized Prompt.

\item AI-Assisted Workflow.

\end{itemize}


The comparison metrics include:


\begin{itemize}

\item Average correctness score.

\item Average runtime.

\item Complexity score.

\end{itemize}



The results will determine whether optimized prompting provides measurable
advantages over direct LLM interaction.



\subsection{Difficulty-Based Analysis}


The benchmark problems will be grouped by difficulty level:


\begin{itemize}

\item Easy.

\item Medium.

\item Hard.

\end{itemize}


This analysis investigates whether prompt engineering provides greater benefits
for more complex algorithmic problems.



\subsection{Algorithm Category Analysis}


Performance differences will also be analyzed across algorithm categories:


\begin{itemize}

\item Array problems.

\item Hash table problems.

\item Tree problems.

\item Graph problems.

\item Dynamic programming problems.

\end{itemize}


This analysis evaluates whether certain algorithm patterns are more suitable
for AI-assisted problem solving.



\subsection{Error Analysis}


Incorrect solutions will be manually analyzed to identify common failure
patterns.


Potential error categories include:


\begin{itemize}

\item Incorrect algorithm selection.

\item Complexity inefficiency.

\item Missing edge cases.

\item Implementation errors.

\end{itemize}



Understanding failure patterns provides insights into the limitations of
LLM-assisted programming.



\subsection{Visualization}


The final results will be presented using:


\begin{itemize}

\item Performance comparison tables.

\item Accuracy comparison charts.

\item Difficulty-based analysis graphs.

\item Algorithm category comparisons.

\end{itemize}


These visualizations will support the evaluation of the research hypotheses.
